% GENERATED FILE. Edit canonical lessons, then run npm run generate:books.
% canonical-source-hash: fnv1a64:d893981800a72204
% canonical-lessons: MR-C429-lilav, MR-C429-taran, MR-C429-gahan-thevne, MR-C429-vasul-karne, MR-C429-bhadepatta

\chapter{Auctions and Pledges}
\label{ch:mr-auctions-and-pledges}
\hlchaptermodality{429}

\begin{chapteropening}
\textbf{By the end of this chapter:} \emph{I can talk about auctions, collateral, pawning, leases and getting money back.}

Complete the last lesson of chapter 429: I can talk about auctions, collateral, pawning, leases and getting money back.
\end{chapteropening}

\section[\texorpdfstring{\mr{लिलाव} (lilāv)}{लिलाव (lilāv)}]{\mr{लिलाव} (lilāv) — an auction}

\label{lesson:MR-C429-lilav}



\begin{quote}
Before the new one: say the Marathi for what one can afford, then the Marathi for economical.
\end{quote}

\subsection*{You'll want to know: \mr{लिलाव}}
\textbf{\mr{लिलाव}} — \emph{lilāv} — \textquotedblleft{}an auction\textquotedblright{}.

\begin{quote}
\mr{जुन्या} \mr{गाड्यांचा} \mr{लिलाव} \mr{उद्या} \mr{आहे}.
\end{quote}

\begin{quote}
\emph{junyā gāḍyāñcā lilāv udyā āhe.}
\end{quote}

\begin{quote}
The auction of old cars is tomorrow.
\end{quote}

Masculine.

\mr{लिलाव} \mr{करणे}, to auction; \mr{लिलावात} \mr{घेणे}, to buy at auction.

\begin{cousinweb}
\textbf{\mr{लिलाव}} (\emph{lilāv}) is a loan from Portuguese \emph{leilão}, \textquotedblleft{}an auction\textquotedblright{}, from the trade of the west coast.
\end{cousinweb}

\begin{grammarlens}[title={What you've built}]
One more word for getting around.
\end{grammarlens}

\subsection*{Your turn}
\begin{itemize}
\raggedright
  \item \emph{Say it:} \emph{lilāv}
  \item \emph{Say it:} \emph{lilāv}, once more
  \item \emph{Say it:} \emph{junyā gāḍyāñcā lilāv udyā āhe.} — The auction of old cars is tomorrow.
  \item \emph{Recall:} say the Marathi for what one can afford, then the Marathi for economical, then say \emph{lilāv} again
\end{itemize}

\begin{tcolorbox}[breakable,colback=teal!4,colframe=teal!35!black,title={Before you move on}]
What does \mr{लिलाव} mean? (\textquotedblleft{}An auction\textquotedblright{}.) Say it once more.
\end{tcolorbox}
\section[\texorpdfstring{\mr{तारण} (tāraṇ)}{तारण (tāraṇ)}]{\mr{तारण} (tāraṇ) — collateral; security}

\label{lesson:MR-C429-taran}



\begin{quote}
Before the new one: say the Marathi for economical, then the Marathi for an auction.
\end{quote}

\subsection*{You'll want to know: \mr{तारण}}
\textbf{\mr{तारण}} — \emph{tāraṇ} — \textquotedblleft{}collateral; security\textquotedblright{}.

\begin{quote}
\mr{त्याने} \mr{कर्जासाठी} \mr{शेत} \mr{तारण} \mr{ठेवलं}.
\end{quote}

\begin{quote}
\emph{tyāne karjāsāṭhī śet tāraṇ ṭhevlaṁ.}
\end{quote}

\begin{quote}
He put up his field as security for the loan.
\end{quote}

Neuter.

\mr{तारण} \mr{ठेवणे}, to give as security for a loan.

\begin{cousinweb}
\textbf{\mr{तारण}} (\emph{tāraṇ}) is inherited: CDIAL lists Marathi \emph{tāraṇ}, \textquotedblleft{}deposit made against a debt\textquotedblright{}, under Sanskrit \emph{tāraṇa}, \textquotedblleft{}crossing, safe passage\textquotedblright{} (CDIAL 5795).
\end{cousinweb}

\begin{grammarlens}[title={What you've built}]
One more word for getting around.
\end{grammarlens}

\subsection*{Your turn}
\begin{itemize}
\raggedright
  \item \emph{Say it:} \emph{tāraṇ}
  \item \emph{Say it:} \emph{tāraṇ}, once more
  \item \emph{Say it:} \emph{tyāne karjāsāṭhī śet tāraṇ ṭhevlaṁ.} — He put up his field as security for the loan.
  \item \emph{Recall:} say the Marathi for economical, then the Marathi for an auction, then say \emph{tāraṇ} again
\end{itemize}

\begin{tcolorbox}[breakable,colback=teal!4,colframe=teal!35!black,title={Before you move on}]
What does \mr{तारण} mean? (\textquotedblleft{}Collateral; security\textquotedblright{}.) Say it once more.
\end{tcolorbox}
\section[\texorpdfstring{\mr{गहाण} \mr{ठेवणे} (gahāṇ ṭhevṇe)}{गहाण ठेवणे (gahāṇ ṭhevṇe)}]{\mr{गहाण} \mr{ठेवणे} (gahāṇ ṭhevṇe) — to pawn}

\label{lesson:MR-C429-gahan-thevne}



\begin{quote}
Before the new one: say the Marathi for an auction, then the Marathi for collateral.
\end{quote}

\subsection*{You'll want to know: \mr{गहाण} \mr{ठेवणे}}
\textbf{\mr{गहाण} \mr{ठेवणे}} — \emph{gahāṇ ṭhevṇe} — \textquotedblleft{}to pawn\textquotedblright{}.

\begin{quote}
\mr{तिने} \mr{सोन्याची} \mr{बांगडी} \mr{गहाण} \mr{ठेवली}.
\end{quote}

\begin{quote}
\emph{tine sonyācī bāṅgḍī gahāṇ ṭhevlī.}
\end{quote}

\begin{quote}
She pawned her gold bangle.
\end{quote}

Past, with -\mr{ने} (he …): \mr{त्याने} \mr{गहाण} \mr{ठेवलं}, \emph{tyāne gahāṇ ṭhevlaṁ}.

The verb agrees with the thing pawned: \mr{बांगडी} \mr{ठेवली}, \mr{घर} \mr{ठेवलं}.

\begin{cousinweb}
\textbf{\mr{गहाण} \mr{ठेवणे}} (\emph{gahāṇ ṭhevṇe}) joins \textbf{\mr{गहाण}}, \textquotedblleft{}a pledge, a pawn\textquotedblright{}, inherited from Sanskrit \emph{grahaṇa}, \textquotedblleft{}seizing, holding\textquotedblright{}, where CDIAL lists Marathi \emph{gahāṇ} (CDIAL 4364), to \textbf{\mr{ठेवणे}}, \textquotedblleft{}to put\textquotedblright{}.
\end{cousinweb}

\begin{grammarlens}[title={What you've built}]
One more word for getting around.
\end{grammarlens}

\subsection*{Your turn}
\begin{itemize}
\raggedright
  \item \emph{Say it:} \emph{gahāṇ ṭhevṇe}
  \item \emph{Say it:} \emph{gahāṇ ṭhevṇe}, once more
  \item \emph{Say it:} \emph{tine sonyācī bāṅgḍī gahāṇ ṭhevlī.} — She pawned her gold bangle.
  \item \emph{Recall:} say the Marathi for an auction, then the Marathi for collateral, then say \emph{gahāṇ ṭhevṇe} again
\end{itemize}

\begin{tcolorbox}[breakable,colback=teal!4,colframe=teal!35!black,title={Before you move on}]
What does \mr{गहाण} \mr{ठेवणे} mean? (\textquotedblleft{}To pawn\textquotedblright{}.) Say it once more.
\end{tcolorbox}
\section[\texorpdfstring{\mr{वसूल} \mr{करणे} (vasūl karṇe)}{वसूल करणे (vasūl karṇe)}]{\mr{वसूल} \mr{करणे} (vasūl karṇe) — to recover money}

\label{lesson:MR-C429-vasul-karne}



\begin{quote}
Before the new one: say the Marathi for collateral, then the Marathi for to pawn.
\end{quote}

\subsection*{You'll want to know: \mr{वसूल} \mr{करणे}}
\textbf{\mr{वसूल} \mr{करणे}} — \emph{vasūl karṇe} — \textquotedblleft{}to recover money\textquotedblright{}.

\begin{quote}
\mr{कंपनीने} \mr{थकबाकी} \mr{वसूल} \mr{केली}.
\end{quote}

\begin{quote}
\emph{kampanīne thakbākī vasūl kelī.}
\end{quote}

\begin{quote}
The company recovered the arrears.
\end{quote}

Past, with -\mr{ने} (he …): \mr{त्याने} \mr{वसूल} \mr{केलं}, \emph{tyāne vasūl kelaṁ}.

The money recovered is the object; \mr{वसूल} \mr{होणे}, to be recovered.

\begin{cousinweb}
\textbf{\mr{वसूल} \mr{करणे}} (\emph{vasūl karṇe}) joins \textbf{\mr{वसूल}}, \textquotedblleft{}recovered, collected\textquotedblright{}, a loan through Persian, usually traced to Arabic \emph{wuṣūl}, \textquotedblleft{}arrival, receipt\textquotedblright{}, to \textbf{\mr{करणे}}.
\end{cousinweb}

\begin{grammarlens}[title={What you've built}]
One more word for getting around.
\end{grammarlens}

\subsection*{Your turn}
\begin{itemize}
\raggedright
  \item \emph{Say it:} \emph{vasūl karṇe}
  \item \emph{Say it:} \emph{vasūl karṇe}, once more
  \item \emph{Say it:} \emph{kampanīne thakbākī vasūl kelī.} — The company recovered the arrears.
  \item \emph{Recall:} say the Marathi for collateral, then the Marathi for to pawn, then say \emph{vasūl karṇe} again
\end{itemize}

\begin{tcolorbox}[breakable,colback=teal!4,colframe=teal!35!black,title={Before you move on}]
What does \mr{वसूल} \mr{करणे} mean? (\textquotedblleft{}To recover money\textquotedblright{}.) Say it once more.
\end{tcolorbox}
\section[\texorpdfstring{\mr{भाडेपट्टा} (bhāḍepaṭṭā)}{भाडेपट्टा (bhāḍepaṭṭā)}]{\mr{भाडेपट्टा} (bhāḍepaṭṭā) — a lease}

\label{lesson:MR-C429-bhadepatta}



\begin{quote}
Before the new one: say the Marathi for to pawn, then the Marathi for to recover money.
\end{quote}

\subsection*{You'll want to know: \mr{भाडेपट्टा}}
\textbf{\mr{भाडेपट्टा}} — \emph{bhāḍepaṭṭā} — \textquotedblleft{}a lease\textquotedblright{}.

\begin{quote}
\mr{त्यांनी} \mr{जमीन} \mr{दहा} \mr{वर्षांच्या} \mr{भाडेपट्ट्याने} \mr{दिली}.
\end{quote}

\begin{quote}
\emph{tyānnī jamīn dahā varṣāñcyā bhāḍepaṭṭyāne dilī.}
\end{quote}

\begin{quote}
They gave the land on a ten-year lease.
\end{quote}

Masculine.

\mr{भाडेपट्ट्याने}, on lease; \mr{पट्टा} here is a written deed.

\begin{cousinweb}
\textbf{\mr{भाडेपट्टा}} (\emph{bhāḍepaṭṭā}) is built from \textbf{\mr{भाडे}}, \textquotedblleft{}rent\textquotedblright{}, and \emph{paṭṭā}, \textquotedblleft{}a strip; a written deed\textquotedblright{}; the origin of \mr{भाडे} is not settled.
\end{cousinweb}

\begin{grammarlens}[title={What you've built}]
One more word for getting around.
\end{grammarlens}

\subsection*{Your turn}
\begin{itemize}
\raggedright
  \item \emph{Say it:} \emph{bhāḍepaṭṭā}
  \item \emph{Say it:} \emph{bhāḍepaṭṭā}, once more
  \item \emph{Say it:} \emph{tyānnī jamīn dahā varṣāñcyā bhāḍepaṭṭyāne dilī.} — They gave the land on a ten-year lease.
  \item \emph{Recall:} say the Marathi for to pawn, then the Marathi for to recover money, then say \emph{bhāḍepaṭṭā} again
\end{itemize}

\begin{tcolorbox}[breakable,colback=teal!4,colframe=teal!35!black,title={Before you move on}]
What does \mr{भाडेपट्टा} mean? (\textquotedblleft{}A lease\textquotedblright{}.) Say it once more.
\end{tcolorbox}
